\scanpage{92}\mspage{45}
c) d'expliciter la construction des $M$
eux-mêmes \unsure{dont|\deleted{dont}} provient la structure
combinatoire envisagée

d) d'expliciter les autres structures qu'on
\unsure{trouve|tire|trouvent} \unsure{à} partir de la donnée combinatoire, telles : la structure de \deleted{$\Lambda^*$}
torseur sous $Z'_\Lambda=\Lambda^*/\pm1$ sur l'ens $\pi$,
le $\Lambda^*$-torseur $\overrightarrow\pi$ \inserted{qui relève $\pi$},
\deleted{l'application} (on
aura $\overrightarrow\pi=\coprod_{i\in\pi}\operatorname{Or}(P_{M,i})$ -- mais il faut
définir l'action de $\Lambda^*$ sur ce $\overrightarrow\pi$)
l'application $\overrightarrow\tau_F:\overrightarrow F_M\to\overrightarrow\pi$, l'opération
de $Z'_\Lambda$ sur $S_M$ \ldots

e) Expliciter comment l'ens des \deleted{\unsure{P}}structures
polyédrales $P_{M,i}$ se déduit de l'une
d'entre elles

f) Examiner, \inserted{pour $\Lambda$ quelconque,} les structures combinatoires
supplémentaires sur les $P_{M,i}$ ou $P_M$,
impliquant $\Lambda$ (p.~ex. l'opération de
$Z'_\Lambda$ sur $S_M$) qui permettent d'énoncer un théorème analogue
au précédent, et de réconstituer
$M$ à symétrie $-1$ près \ldots
% archive sheet number: 91
